\subsection{Quadratic Forms and Polar Forms}

\begin{definition} \label{def:quadratic-form-and-polar-form} \hlt{Quadratic Forms and Polar Forms}\\
A \hlt{quadratic form} on $\mathcal{V}$ is a function $Q:\mathcal{V}\to\mathbb{F}$ such that:
\begin{enumerate}[label=(\roman*)]
\setlength{\itemsep}{0pt}
\item $Q(a\mathbf{v})=a^2Q(\mathbf{v})$ for all $a\in\mathbb{F}$ and $\mathbf{v}\in\mathcal{V}$;
\item the \hlt{polar form} $B_Q(\mathbf{u},\mathbf{v})\defeq Q(\mathbf{u}+\mathbf{v})-Q(\mathbf{u})-Q(\mathbf{v})$ is bilinear.
\end{enumerate}
\end{definition}

\begin{proposition} \label{prop:quadratic-forms-symmetric-bilinear-forms} \hlt{Quadratic Forms and Symmetric Bilinear Forms}\\
For every quadratic form $Q$, its polar form $B_Q$ is symmetric.\\
If $\operatorname{char}(\mathbb{F})\neq2$, maps $\{B:\operatorname{Sym}(B)\} \underset{Q\mapsto B_Q}{\overset{B\mapsto Q_B}{\rightleftarrows}} \{Q:Q\text{ quadratic}\}$ and $Q_B(\mathbf{x})\defeq\frac{1}{2}B(\mathbf{x},\mathbf{x})$ are inverse correspondences between symmetric bilinear forms and quadratic forms.
\begin{proof} First, $B_Q(\mathbf{u},\mathbf{v}) =Q(\mathbf{u}+\mathbf{v})-Q(\mathbf{u})-Q(\mathbf{v}) =B_Q(\mathbf{v},\mathbf{u})$, so $B_Q$ is symmetric.
\begin{enumerate}[label=\roman*.]
\setlength{\itemsep}{0pt}
\item[$(\Rightarrow)$] If $\operatorname{Sym}(B)$, then $Q_B(a\mathbf{x})=a^2Q_B(\mathbf{x})$ and $B_{Q_B}(\mathbf{u},\mathbf{v}) =\frac{1}{2}\bigl(B(\mathbf{u}+\mathbf{v},\mathbf{u}+\mathbf{v})-B(\mathbf{u},\mathbf{u})-B(\mathbf{v},\mathbf{v})\bigr) =B(\mathbf{u},\mathbf{v})$.
\item[$(\Leftarrow)$] If $Q$ is quadratic, then $Q_{B_Q}(\mathbf{x}) =\frac{1}{2}B_Q(\mathbf{x},\mathbf{x}) =\frac{1}{2}\bigl(Q(2\mathbf{x})-2Q(\mathbf{x})\bigr) =Q(\mathbf{x})$.
\end{enumerate}
\hlt{Structural Commentary}: Away from characteristic $2$, quadratic and symmetric bilinear data coincide.\\
\hlt{Key Insight}: Dividing by $2$ recovers diagonal data from polar data.
\end{proof}
\end{proposition}

\begin{proposition} \label{prop:polar-forms-alternating-in-characteristic-two} \hlt{Polar Forms in Characteristic Two}\\
If $\operatorname{char}(\mathbb{F})=2$, then $\operatorname{Alt}(B_Q)$ for every quadratic form $Q$ on $\mathcal{V}$.
\begin{proof} For $\mathbf{x}\in\mathcal{V}$, $B_Q(\mathbf{x},\mathbf{x}) =Q(2\mathbf{x})-Q(\mathbf{x})-Q(\mathbf{x}) =Q(\mathbf{0})=0$. Thus $\operatorname{Alt}(B_Q)$.\\
\hlt{Structural Commentary}: In characteristic $2$, every polar form is alternating.\\
\hlt{Key Insight}: The diagonal of the polar form vanishes automatically because both $2\mathbf{x}$ and $2Q(\mathbf{x})$ vanish.
\end{proof}
\end{proposition}

\begin{example} \label{ex:characteristic-two-polar-form-loses-quadratic-data} \hlt{Characteristic Two Loses Diagonal Data}\\
Assume $\operatorname{char}(\mathbb{F})=2$. On $\mathbb{F}^2$, let $Q_0(x,y)=xy$ and $Q_1(x,y)=xy+x^2$.\\
Then $B_{Q_0}((x,y),(u,v))=xv+uy=B_{Q_1}((x,y),(u,v))$, and $Q_0(1,0)=0\neq1=Q_1(1,0)$.\\
Indeed, the added square has zero polar contribution, where $(x+u)^2-x^2-u^2=2xu=0$.\\
More generally, for any linear $L:\mathcal{V}\to\mathbb{F}$, $Q$ and $Q+L^2$ have the same polar form.\\
\hlt{Structural Commentary}: In characteristic $2$, squares are invisible to polarization.\\
\hlt{Key Insight}: Characteristic $2$ quadratic forms need data beyond their polar forms.
\end{example}

\begin{example} \label{ex:characteristic-two-alternating-versus-nonalternating-symmetric} \hlt{Alternating versus Nonalternating Symmetric Forms}\\
Assume $\operatorname{char}(\mathbb{F})=2$.\\
The symmetric matrix $H=\begin{bmatrix}0&1\\1&0\end{bmatrix}$ defines a nonsingular alternating form on $\mathbb{F}^2$, since $\begin{bmatrix}x&y\end{bmatrix}H\begin{bmatrix}x\\y\end{bmatrix}=2xy=0$.\\
Thus $H$ is not congruent to any nonzero diagonal form: alternating means all self-pairings vanish, while a nonzero diagonal form has some basis vector with nonzero self-pairing. By contrast, $[1]$ is symmetric and nonalternating.\\
\hlt{Structural Commentary}: In characteristic $2$, symmetric forms split into alternating and nonalternating parts.\\
\hlt{Key Insight}: Nonsingular alternating blocks are invisible on the diagonal.
\end{example}
